\chapter{Il confronto finale con Elasticsearch}\label{chapter:confronto}

\section{Le configurazioni a confronto}\label{sec:configurazioni}

\todo[inline]{Bozza del 28/09, prima del terzo lotto: struttura e ragionamento
fissati, i numeri (\dacompletare) arrivano da
\texttt{2026-09-25\_known-item-umane} a raccolta chiusa.}

Il confronto finale si regge sulle known-item umane
(sezione~\ref{sec:query}): \dacompletare{} query scritte da persone che non
conoscono i motori, una per atto, su \dacompletare{} atti estratti con un seme,
metà per fonte. Per ogni query c'è un atto giusto solo, quindi la misura
principale è l'MRR@10, accanto a tre quantità che un utente sente di più:
l'atto al primo posto, l'atto fra i primi dieci, la ricerca a vuoto. I bisogni
aperti delle 24 query restano un confronto secondario: i giudizi del modello
concordano con quelli umani con un kappa pesato di 0,370
(sezione~\ref{sec:giudizi}), troppo poco per reggere un capitolo, e il resto
del pool non ha un giudizio umano.

Le configurazioni sono di due tipi, e si misurano in due modi.

\paragraph{Dall'app.} Due, sulle stesse query passate da
\texttt{app:eval-run-queries}, cioè dalla ricerca vera di Documentale:
\begin{itemize}
  \item \textbf{ES}: Elasticsearch come lo usa Documentale in produzione, con il
        \texttt{multi\_match} \texttt{best\_fields} e l'operatore
        \texttt{and} (capitolo~\ref{chapter:sistemi});
  \item \textbf{KC}: Koskidex innestato con il profilo consigliato
        (sezione~\ref{sec:cosa-cambia}): parole libere di stare in campi
        diversi, numeri esatti, elisione; senza vettori.
\end{itemize}

\paragraph{Koskidex piatto.} Sei, con \texttt{scripts/evaluate} sul corpus a
schede che esce dall'esportazione di Documentale, per vedere dove porterebbero
le correzioni dei capitoli~\ref{chapter:lessicale} e~\ref{chapter:ibrido} se
entrassero nell'app: \textbf{K0}, il punto di partenza (congiuntivo,
punteggio euristico); \textbf{LA} e \textbf{LT}, BM25 con la frequenza
mescolata in recupero disgiuntivo e congiuntivo; \textbf{A} e \textbf{B},
le due configurazioni ibride del capitolo~\ref{chapter:ibrido}, quella per le
frasi e quella per gli identificativi; \textbf{V}, i soli vettori. Le
configurazioni piatte non passano dall'app e non hanno i refusi, quindi fra
loro si confrontano alla pari, con ES e KC solo nell'ordine di grandezza.

\paragraph{Il tetto del riordino.} Un'ultima coppia non è una proposta:
\textbf{RR-LA} e \textbf{RR-A} sono i primi cento di LA e di A riordinati da un
cross-encoder multilingue (\texttt{bge-reranker-v2-m3}), che legge query e
documento insieme e costa secondi a query su questo hardware
(\texttt{2026-09-28\_reranker}). Dicono quanto vale, al massimo, riordinare
quello che il lessicale trova.

Ogni metrica si riporta su tutte le query, per fonte (Crispiano, dove chi
scriveva aveva letto il testo intero e l'indice ha la scheda, e Friuli Venezia
Giulia, dove l'atto è la sola scheda) e per query con e senza un numero. Le
sette previsioni sono state scritte il 25 settembre, prima che arrivasse una
sola risposta.

\section{Pertinenza, per famiglia di query e per fonte}\label{sec:pertinenza}

\begin{table}[ht]
  \centering
  \small
  \begin{tabular}{lrrrrr}
    \toprule
    & tutte & Crispiano & Friuli V.G. & con numero & senza numero \\
    \midrule
    ES & \dacompletare & \dacompletare & \dacompletare & \dacompletare & \dacompletare \\
    KC & \dacompletare & \dacompletare & \dacompletare & \dacompletare & \dacompletare \\
    \midrule
    K0 & \dacompletare & \dacompletare & \dacompletare & \dacompletare & \dacompletare \\
    LA & \dacompletare & \dacompletare & \dacompletare & \dacompletare & \dacompletare \\
    LT & \dacompletare & \dacompletare & \dacompletare & \dacompletare & \dacompletare \\
    A & \dacompletare & \dacompletare & \dacompletare & \dacompletare & \dacompletare \\
    B & \dacompletare & \dacompletare & \dacompletare & \dacompletare & \dacompletare \\
    V & \dacompletare & \dacompletare & \dacompletare & \dacompletare & \dacompletare \\
    \midrule
    RR-LA & \dacompletare & \dacompletare & \dacompletare & \dacompletare & \dacompletare \\
    RR-A & \dacompletare & \dacompletare & \dacompletare & \dacompletare & \dacompletare \\
    \bottomrule
  \end{tabular}
  \caption{Known-item umane: MRR@10 e, fra parentesi, la quota di ricerche a
  vuoto. ES e KC dall'app, le altre con Koskidex piatto sul corpus a schede.}
  \label{tab:known-item-umane}
\end{table}

\paragraph{Che cosa scrive chi cerca.} \dacompletare{} query su
\dacompletare{} contengono un numero; avevo previsto meno di una su quattro,
perché chi ricorda un atto dopo settimane ne ricorda l'argomento, non il
protocollo. \dacompletare{} persone su \dacompletare{} hanno risposto «non
saprei» per almeno un atto.
\todo[inline]{Previsione 1. Se le query con un numero sono poche, le colonne
«con numero» vanno lette come aneddoti: dirlo con il loro conteggio.}

\paragraph{Il campo unico anche sulle ricerche per contenuto.} Il difetto della
sezione~\ref{sec:numero-comune} era stato trovato sulle ricerche per numero e
comune. Qui la domanda è se colpisca anche chi cerca per argomento: con
l'operatore \texttt{and} di produzione tutte le parole devono stare nello
stesso campo della scheda, e chi scrive aggiunge facilmente il comune, un anno o
una parola del testo. ES va a vuoto nel \dacompletare{} delle query e trova
l'atto fra i primi dieci nel \dacompletare; KC, con le parole libere di stare
in campi diversi e l'elisione, va a vuoto nel \dacompletare{} e fa
\dacompletare{} di MRR@10 sopra ES.
\todo[inline]{Previsioni 2 e 3: più del 30\% a vuoto per ES, KC almeno 0,10
sopra con almeno 10 punti di vuoti in meno. Se tengono tutte e due, il
difetto del campo unico riguarda anche le ricerche per contenuto: è la frase
da scrivere.}

\paragraph{Congiuntivo o disgiuntivo.} Una parola in più, o scritta
diversamente da come è nell'atto, svuota il recupero congiuntivo. Sulle query
senza numero LA fa \dacompletare{} e LT \dacompletare; su quelle con un numero
\dacompletare.
\todo[inline]{Previsione 4: LA almeno 0,05 sopra LT senza numero, LT sopra LA
con il numero. Collegare alla sezione~\ref{sec:scelta}: la regola per tipo di
query e il classificatore sulle stesse query (\texttt{scelta-umane}).}

\paragraph{I vettori.} A, che aggiunge ai candidati di LA i cento più vicini,
fa \dacompletare{} sulle query senza numero, \dacompletare{} sopra LA; i soli
vettori (V) \dacompletare, e sulle query con un numero \dacompletare.
\todo[inline]{Previsioni 5 e 7: A almeno 0,03 sopra LA senza numero; V entro
0,10 da LA senza numero e sotto 0,30 con il numero. È la verifica, su query
vere, di quello che il capitolo~\ref{chapter:ibrido} aveva visto sulle
collezioni pubbliche.}

\paragraph{Le due fonti.} Su Crispiano chi scriveva aveva letto il testo
intero, e l'indice ha la scheda: le parole che ricorda possono non esserci. LA
fa \dacompletare{} su Crispiano e \dacompletare{} sul Friuli Venezia Giulia, KC
\dacompletare{} e \dacompletare.
\todo[inline]{Previsione 6: almeno 0,10 sotto su Crispiano per LA e KC.
Rimandare alla sezione~\ref{sec:scheda-testo} per scheda contro testo intero
sui soli atti di Crispiano, e alla sezione~\ref{sec:italiana} per stopword e
stemmer italiani sulle stesse query.}

\paragraph{Quanto vale riordinare.} Il cross-encoder porta LA da
\dacompletare{} a \dacompletare{} di MRR@10 e A da \dacompletare{} a
\dacompletare; degli atti che il primo stadio aveva fra l'undicesimo e il
centesimo posto ne porta fra i primi dieci \dacompletare{} su \dacompletare.
Sulle collezioni pubbliche il guadagno è più piccolo e ha un tetto: SciFact
passa da 0,676 a 0,723 di nDCG@10, NFCorpus da 0,306 a 0,331, e su NFCorpus
il primo stadio porta fra i cento soltanto il 24\% dei documenti giusti, gli
unici che il riordino può vedere. Sulle ricerche \texttt{<numero> <comune>}
porta LA da 0,833 a 0,953, contro 0,960 del recupero congiuntivo: degli atti
fra l'undicesimo e il centesimo posto ne porta fra i primi dieci 12 su 13.
Avevo previsto che restasse sotto LT perché un cross-encoder non pesa un
numero esatto più di uno simile; resta sotto, ma di 0,007, e il ragionamento
era sbagliato. Il prezzo è di 5 secondi di mediana a query sulle schede e di
12 sugli abstract di SciFact, contro il millisecondo di Koskidex: la misura di
un tetto, non una configurazione da mettere in produzione.
\todo[inline]{Previsioni umane di \texttt{reranker}: almeno +0,05 su LA e
+0,03 su A, almeno metà degli atti dall'11° al 100° posto portati fra i primi
dieci.}

\section{Latenza, memoria, dimensione dell'indice}\label{sec:latenza}

La sezione~\ref{sec:prestazioni-innesto} ha dato i tempi dei due motori
passando dall'app, con Koskidex nativo ed Elasticsearch nella macchina
virtuale: un ordine di grandezza, non un confronto alla pari. Qui i due motori
girano nella stessa macchina virtuale, con le stesse 4 CPU e la stessa rete:
Koskidex è lo stesso binario compilato per Linux, in un container accanto a
quello di Elasticsearch, con l'indice riempito dall'app nella configurazione
consigliata. Le richieste sono quelle che manda l'app, copiate dal codice: per
Elasticsearch il \texttt{multi\_match} di produzione con i sei campi pesati,
che restituisce anche il testo di ogni risultato; per Koskidex la ricerca con
i soli identificativi. Le query sono 400: le 300 known-item automatiche, le 24
del confronto e le 76 known-item umane raccolte fino al 28 settembre. Tutto
sta in \texttt{risultati/esperimenti/2026-09-28\_carico/}, previsioni
comprese.

Koskidex tiene in memoria le ultime 1.024 risposte, e non c'è un'impostazione
per spegnere la cache: ripetere le stesse query misurerebbe la cache. Ogni
richiesta porta quindi un numero diverso di spazi in coda alla query, che il
tokenizer ignora e la chiave della cache no; lo strumento verifica, prima di
misurare, che con gli spazi nessuna delle 400 risposte cambi, su nessuno dei
due motori.

\paragraph{Una ricerca alla volta.} Ogni query cinque volte, in un ordine
mescolato con un seme fisso. Il p50 di Koskidex è un terzo di quello di
Elasticsearch, e l'avevo previsto; avevo previsto anche che la differenza
stesse soprattutto nel peso della risposta, perché Elasticsearch restituisce il
testo degli atti che l'app poi scarta. Non è così: senza il testo Elasticsearch
non va più veloce (6,77~ms), perché con l'operatore \texttt{and} le risposte
sono piccole, 160 byte di mediana. La differenza sta nel motore. La coda invece
era quasi la stessa (tabella~\ref{tab:carico-sequenziale}, colonna
\emph{prima}): 68 query su 400 superavano gli 8~ms, ed erano le query lunghe,
piene di parole comuni.

\begin{table}[ht]
  \centering
  \begin{tabular}{lrrr}
    \toprule
    & p50 & p95 & p99 \\
    \midrule
    Elasticsearch, come l'app & 6,32 & 17,66 & 24,26 \\
    Elasticsearch, senza il testo & 6,77 & 19,60 & 26,09 \\
    Koskidex, prima & 2,12 & 12,36 & 18,31 \\
    Koskidex, dopo & \textbf{1,12} & \textbf{2,74} & \textbf{4,79} \\
    \bottomrule
  \end{tabular}
  \caption{Una ricerca alla volta, in millisecondi, 400 query per cinque
  passate, i due motori nella stessa macchina virtuale. \emph{Prima} è
  Koskidex \texttt{ffa38ab}, \emph{dopo} \texttt{510b9d2}, con le due correzioni
  di questa sezione.}
  \label{tab:carico-sequenziale}
\end{table}

\paragraph{Sotto carico.} Da 1 a 32 client che mandano ricerche senza pause,
20 secondi per livello. Avevo previsto che Koskidex reggesse almeno una volta e
mezza le ricerche al secondo di Elasticsearch. Ne reggeva 1,14 volte: 653
contro 571 (tabella~\ref{tab:carico}). Con un client solo usava già il 116\% di
CPU, più di un core per una ricerca alla volta, e a 32 client il suo p99 era il
doppio di quello di Elasticsearch.

\begin{table}[ht]
  \centering
  \begin{tabular}{rrrrrrr}
    \toprule
    & \multicolumn{2}{c}{Elasticsearch} & \multicolumn{2}{c}{Koskidex, prima} &
    \multicolumn{2}{c}{Koskidex, dopo} \\
    \cmidrule(lr){2-3} \cmidrule(lr){4-5} \cmidrule(lr){6-7}
    client & ricerche/s & p99 & ricerche/s & p99 & ricerche/s & p99 \\
    \midrule
    1 & 134,6 & 22,8 & 233,3 & 19,2 & 705,9 & 4,9 \\
    4 & 471,7 & 26,6 & 588,9 & 29,1 & 1.707,9 & 7,9 \\
    8 & 571,1 & 46,5 & 653,0 & 48,5 & 2.313,8 & 13,9 \\
    16 & 550,4 & 69,6 & 622,6 & 105,0 & 2.588,7 & 25,3 \\
    32 & 559,4 & 103,1 & 647,2 & 214,6 & \textbf{2.614,0} & \textbf{42,8} \\
    \bottomrule
  \end{tabular}
  \caption{Sotto carico, stesse 4 CPU: ricerche al secondo e p99 in
  millisecondi. Elasticsearch e Koskidex prima saturano le CPU dagli 8 client
  in su (380-390\%); Koskidex dopo arriva al 319\% a 8 client e al 360\% a
  32, e il resto va alla rete.}
  \label{tab:carico}
\end{table}

\paragraph{Dove spendeva Koskidex.} Il profilo a 8 client ha dato la risposta, e
non era quella che avevo previsto: avevo scommesso sull'espansione dei refusi e
sulla scrittura del JSON della risposta. Koskidex allocava 2,57~MB a ricerca, e
quella memoria si pagava in CPU: il 20\% dei campioni era nel lock
dell'allocatore di Go e un altro 7\% nel garbage collector. Il 43\% delle
allocazioni veniva dalla funzione che, per ogni parola della query, costruisce
una voce per ogni documento che la contiene, anche per \emph{di} o \emph{per},
e solo dopo scarta i documenti a cui manca una parola; un altro 40\% dalla
distanza fra parole, che allocava una matrice intera a ogni confronto, e dalla
raccolta dei candidati con refuso.

\paragraph{Due correzioni che non cambiano un risultato.} A differenza delle
correzioni dei capitoli precedenti non stanno dietro un'impostazione, e il
motivo è la condizione stessa del lavoro: non devono cambiare né l'insieme né l'ordine né un
punteggio, e se ne cambiassero uno non sarebbero queste correzioni. La prima
calcola la distanza con tre righe riusate al posto della matrice, la stessa
ricorrenza, e prende i bigrammi come sottostringhe. La seconda, con il recupero
congiuntivo, trova prima i documenti che contengono tutte le parole, partendo
da quella con meno occorrenze come fa Elasticsearch, e calcola il punteggio
solo per loro, parola per parola nell'ordine della query: le somme in virgola
mobile restano le stesse bit per bit. La prova è stata fatta prima di misurare
i tempi: l'impronta di ogni risposta del server, identificativi e punteggi
byte per byte, per 429 query (le 400 più 29 inventate con gli operatori, parole
di una lettera, numeri, refusi e parole inesistenti), coincide prima e dopo
per tutte, a 10.018 documenti e a 100.000; le metriche di
\texttt{scripts/evaluate}, query per query, coincidono su SciFact, NFCorpus e le
known-item automatiche nelle tre configurazioni di confronto; i test di
Koskidex passano, compreso quello che congela l'ordine di partenza.

Le otto previsioni scritte prima (\texttt{2026-09-28\_prestazioni}) sono tutte
confermate, quasi sempre con molto margine: le soglie erano prudenti. Le
allocazioni scendono da 2,57 a 0,44~MB a ricerca, l'83\% in meno. Una ricerca
alla volta costa 1,12~ms di p50 e 2,74 di p95, e il p95 scende di più del p50:
le correzioni hanno tolto proprio il costo delle query lunghe, che faceva la
coda. Sotto carico Koskidex regge 2.614 ricerche al secondo, 4,6 volte
Elasticsearch sulle stesse CPU, e a 32 client il suo p99 è 43~ms contro 103.
Con un client la CPU scende dal 116\% all'83\%. Nel profilo nuovo il lock
dell'allocatore e il garbage collector escono dalle prime righe, e il 43\% dei
campioni è nelle chiamate di sistema della rete: il motore non è più la parte
lenta della richiesta. Nativo, sugli 8 core del Mac, Koskidex passa da 975 a
4.100 ricerche al secondo.

\paragraph{Una terza correzione, e due previsioni sbagliate.} Delle allocazioni
che restavano, l'86\% veniva da una riga sola: l'insieme che toglie i doppioni
fra le parole del vocabolario che condividono un bigramma con quella cercata,
migliaia per i bigrammi comuni, quasi tutte scartate subito dopo perché troppo
lunghe o troppo corte. La correzione scarta prima e toglie i doppioni solo fra
quelle che restano; il controllo dipende solo dalla parola, quindi le parole
trovate sono le stesse, nello stesso ordine
(\texttt{2026-09-28\_allocazioni}). La prova è la stessa: 429 impronte su 429
a tutte e due le dimensioni, 9 valutazioni su 9, più un test che confronta
42.420 ricerche di parole con il percorso di prima. Misurato contro un
\emph{prima} rifatto nella stessa mezz'ora, le allocazioni scendono da 0,435 a
0,081~MB a ricerca, il p50 di una ricerca alla volta nel container da 1,14 a
0,92~ms, e la capacità nel container da 2.622 a 3.779 ricerche al secondo,
+44\%, 6,6 volte Elasticsearch; nativo da 4.090 a 6.376. Cinque
previsioni su sette. Una soglia era una quota, meno del 10\% dei byte allocati
per quella funzione: è scesa al 30\%, ma perché è sceso tutto il resto con lei,
e in valore assoluto è passata da 385 a 25~KB a ricerca. L'altra si aspettava
un guadagno forte a 100.000 documenti, dove pensavo che le liste dei bigrammi
fossero dieci volte più lunghe; ma la copia ripete gli stessi testi, il
vocabolario non cresce, e la coda lì viene dalle liste dei documenti: p95 e p99
scendono solo del 4 e del 7\%.

\paragraph{Memoria e archivi più grandi.} Per vedere come cresce il costo ho
copiato gli stessi atti fino a 50.000 e 100.000 documenti, con identificativi
diversi e lo stesso testo: dati inventati, che per tempi e memoria vanno bene e
per la pertinenza no (tabella~\ref{tab:crescita}).

\begin{table}[ht]
  \centering
  \small
  \begin{tabular}{lrrrrrr}
    \toprule
    & \multicolumn{2}{c}{memoria} & \multicolumn{2}{c}{indice su disco} &
    \multicolumn{2}{c}{indicizzazione} \\
    \cmidrule(lr){2-3} \cmidrule(lr){4-5} \cmidrule(lr){6-7}
    documenti & ES & Koskidex & ES & Koskidex & ES & Koskidex \\
    \midrule
    10.018 & 1.431 MB & 155 MB & 3,5 MB & 6,3 MB & 1,1-1,7 s & 1,9-2,1 s \\
    50.000 & 1.439 MB & 651 MB & 17,9 MB & 32,8 MB & 2,8 s & 6,2 s \\
    100.000 & 1.433 MB & 1.213 MB & 35,5 MB & 65,6 MB & 5,3 s & 12,9 s \\
    \bottomrule
  \end{tabular}
  \caption{Memoria a riposo dopo un riavvio, indice su disco e tempo di
  indicizzazione, con gli atti copiati. A 10.018 documenti l'indicizzazione è
  quella dall'app.}
  \label{tab:crescita}
\end{table}

Alla dimensione dell'albo Koskidex occupa un decimo della memoria di
Elasticsearch, e sotto carico 272~MB al massimo contro 1.566, 195 dopo la terza
correzione. Ma tiene tutto in
memoria, con gli identificativi dei documenti e i nomi dei campi come stringhe
in ogni occorrenza: circa 12~KB a documento, e a 100.000 documenti la memoria è
quasi quella di Elasticsearch, che ha un heap fisso. Le correzioni non la
toccano (1.213 contro 1.184~MB), perché non toccano le strutture dell'indice.
Toccano invece la coda: a 100.000 documenti il p95 di Koskidex era di 133~ms e
il p99 di 201, dieci volte quelli a 10.018, contro i 21 e 36 di Elasticsearch;
dopo sono 10,4 e 15,5, sotto quelli di Elasticsearch, e 9,9 e 13,6 dopo la
terza.

\paragraph{Un difetto trovato per strada.} Il test che confronta i due percorsi
su query casuali ha trovato una cosa che non cercava, e che ha avuto il suo
esperimento (\texttt{2026-09-28\_ordine-fisso}). Per una parola corta con
un refuso ammesso, quattro lettere o meno con la tolleranza automatica,
Koskidex non usa i bigrammi ma scorre tutto il vocabolario, e il vocabolario è
una mappa di Go, che si percorre in un ordine diverso a ogni chiamata. L'ordine
delle parole evidenziate cambia da una ricerca all'altra; e se un documento
contiene due parole alla stessa distanza da quella cercata, quale delle due gli
viene accreditata, e quindi il suo punteggio BM25, dipende dal caso: in un test
con due documenti che contengono le stesse due parole con frequenze diverse,
i due documenti si scambiano di posto da una chiamata all'altra. Lo stesso
succede con la ricerca per sottostringa, che scorre il vocabolario allo stesso
modo, e dopo un riavvio anche altrove, perché l'istantanea su disco scrive i
documenti in ordine di mappa. La correzione è un'impostazione,
\texttt{StableTermOrder}, spenta per default, che visita i termini trovati in
ordine lessicografico. Sei previsioni su sette hanno tenuto: nessun numero di
\texttt{scripts/evaluate} dipendeva dall'ordine della mappa, perché lo
strumento cerca senza refusi, e cinque ripetizioni non variano mai, spente o
accese; Documentale, con il punteggio euristico, dà le stesse 429 risposte;
il costo sta nel rumore. È caduta quella che dava per uguali le medie con
l'impostazione accesa e spenta: su SciFact con BM25 disgiuntivo passano da
0,6694 a 0,6670. Anche l'ordine di ingresso nell'indice, che le misure della
tesi hanno sempre usato, è arbitrario, e decide a parità di distanza quale
espansione per prefisso viene accreditata; la differenza, due query su 300, è
quanto pesa quella regola (sezione~\ref{sec:espansioni}).

\paragraph{Cosa ne segue per Documentale.} All'albo di un comune, diecimila
atti, Koskidex costa un decimo della memoria di Elasticsearch e risponde in
meno di un quinto del tempo, e sotto carico regge più di sei volte le ricerche
sulle stesse CPU. A centomila documenti il tempo regge ancora, la memoria no: lì
serve un indice più compatto, con identificativi e campi interi, che
ridurrebbe anche l'indice su disco, oggi il doppio di quello di
Elasticsearch, e l'indicizzazione, 2,4 volte più lenta. I limiti delle
misure: il client gira sullo stesso Mac, fuori dalla macchina virtuale, e tiene
le connessioni aperte, mentre l'app in PHP ne apre in genere una per ricerca;
Elasticsearch ha l'heap fissato a 1~GB, e con meno memoria si comporterebbe
diversamente.

\section{Dove un motore piccolo è competitivo, e dove no}\label{sec:competitivo}

\todo[inline]{Bozza del 28/09: le affermazioni sui costi hanno già i loro
numeri (sezione~\ref{sec:latenza}); quelle sulla pertinenza aspettano la
sezione~\ref{sec:pertinenza}.}

La domanda della tesi non era se Koskidex batte Elasticsearch, ma se un motore
piccolo, innestato in un documentale al posto di uno grande, regge, e dove
smette di reggere. Le risposte sono di tre tipi.

\paragraph{Dove regge.} All'albo di un comune, diecimila atti, Koskidex
risponde in 0,92~ms di p50 contro 6,32, regge 6,6 volte le ricerche al secondo
sulle stesse quattro CPU e occupa un decimo della memoria. Sulle ricerche per
numero e comune il profilo consigliato porta l'atto giusto al primo posto nel
91,7\% dei casi, contro l'1,3\% di produzione, e le date si trovano in
qualunque formato. Sulle ricerche per contenuto scritte da persone KC fa
\dacompletare{} di MRR@10 contro \dacompletare{} di ES.
\todo[inline]{Se la previsione 3 cade, questa frase cambia verso: dire allora
che la differenza di pertinenza sta nelle impostazioni di produzione, non nel
motore, e che Elasticsearch con le stesse correzioni farebbe lo stesso
(la parità del capitolo~\ref{chapter:sistemi}).}

\paragraph{Dove il vantaggio non è del motore.} Quasi tutto il guadagno di
pertinenza viene da impostazioni che Elasticsearch ha già, o potrebbe avere:
parole in campi diversi con \texttt{cross\_fields}, l'elisione con il suo
filtro, le date con un \emph{char filter}; i numeri esatti non hanno
un'opzione, ma si ottengono separando nella query i termini numerici e
cercandoli senza fuzziness. Il capitolo~\ref{chapter:sistemi} ha
mostrato che con le stesse impostazioni i due motori trovano gli stessi
insiemi. Il merito di Koskidex, in questa tesi, è stato di rendere visibili e
misurabili i difetti di una configurazione, non di esserne immune.

\paragraph{Dove non regge.} La memoria: Koskidex tiene tutto in memoria, con
identificativi e nomi dei campi come stringhe in ogni occorrenza, e a centomila
documenti ne usa 1.213~MB, quasi quanto l'heap fisso di Elasticsearch. L'indice
su disco è il doppio e l'indicizzazione 2,4 volte più lenta. Mancano le cose che
un documentale grande dà per scontate e che questa tesi non ha misurato:
repliche, shard, un ecosistema di strumenti. E la pertinenza sulle ricerche
aperte resta non misurata con giudizi umani.
\todo[inline]{Chiudere con una frase sola: per chi, con quali numeri di
archivio e di utenti, un motore piccolo è la scelta giusta.}
